\documentclass[11pt]{article}
\usepackage[a4paper,margin=1in]{geometry}
\usepackage[T1]{fontenc}
\usepackage{booktabs}
\usepackage{hyperref}
\usepackage{enumitem}
\setlist{nosep}

\title{Urban Data Integration Platform\\Operations and Maintenance Report}
\author{Data-intensive Computing}
\date{}

\begin{document}
\maketitle

\section{Purpose and scope}

This report adds a second data release to the Urban Data Integration Platform.
The work builds on the completed Bronze, Silver, Gold, and data-product layers.
It does not modify the existing ingestion, integration, or analytical-product
implementation. Instead, the small operations layer in
\texttt{src/operations/pipeline.py} creates and applies a realistic release,
records each operation, and refreshes the outputs affected by new data.

The final release contains later taxi trips, later hourly weather observations,
and later hourly air-quality observations. It also tests two compatible schema
additions and duplicate handling. The workflow is reproducible through four
scripts under \texttt{scripts/operations/}.

\section{Incremental release and insertion}

\subsection{Release design}

The generator writes a separate release below \texttt{data/updates/}; the
original source files in \texttt{data/raw/} remain unchanged. The final release contains 508,675 new taxi trips, 8,490 copied taxi duplicates, 168 weather rows with \texttt{humidity}, and 168 air-quality rows with \texttt{aqi}.


Taxi records are sampled from valid source trips and shifted after the latest
original timestamp while preserving trip duration. The 508,675 new trips are
5.32\% of the 9,554,778-row taxi source, meeting the required 5--10\% range.
The 8,490 unchanged copied trips are 1.64\% of the 517,165-row taxi release,
meeting the required 1--2\% duplicate range. Weather and air quality each add
one week of hourly data. Humidity is generated in the range 20--100; AQI is
generated in the range 0--500.

\subsection{Validation, duplicates, and schema evolution}

The pipeline uses the existing common loading and standardisation helpers.
Before writing, it checks required fields, original dataset validity rules,
taxi-zone references, documented numeric ranges, duplicates, and unexpected
columns. Rejected records receive a \texttt{rejection\_reason} and are written
to \texttt{data/rejected/<dataset>/}; valid records continue through the
pipeline.

Weather and air quality have usable observation keys. Their rows are compared
with existing keys only within the relevant year/month scope. Taxi has no
stable trip identifier, so the pipeline calculates a SHA-256 hash over shared
fields and uses an anti join to select rows absent from Bronze. This approach
is insert-only: an unchanged Bronze record is never overwritten. Re-running
the same release therefore inserts no additional records.

Schema evolution is deliberately narrow. \texttt{humidity} for weather and
\texttt{aqi} for air quality are explicitly allowed. Any other new column is
reported and excluded, while a missing required column prevents unsafe input
from being understood as a valid update.

\section{Maintaining analytical consistency}

The dependent flow is:

\begin{verbatim}
update -> Bronze append -> changed Silver/Gold tables -> affected products
\end{verbatim}

Only datasets with inserted rows are rebuilt. Product dependencies are kept in
one visible map. Daily mobility and taxi-zone statistics depend on taxi trips;
weather impact depends on taxi and weather; air-quality impact depends on taxi
and air quality. The final release changed all three source datasets, so all
four products refreshed.

The product refresh is a full rebuild only when its dependencies change. This
is an appropriate simple design because the products are whole-history
aggregates, and it avoids work after an idempotent repeat. A later improvement
would replace only affected month partitions in daily mobility. Using the two
new schema fields in an analytical product would be a deliberate product
change; existing products remain compatible because they select the columns
they require.

\section{Monitoring and rejected records}

Every update, layer rebuild, and product refresh appends one compact record to
the Delta table \texttt{data/monitoring/operations\_runs}. It stores timestamp,
dataset, stage, processed/inserted/rejected counts, total and validation time,
schema version, validation failures, and schema changes. The last two fields
are JSON text, making new validation rules possible without changing the
monitoring-table schema.

The monitoring script runs four Spark SQL summaries: rejection totals by
dataset, average execution time, rejected rows and reasons for each execution,
and timing change between runs. The final run showed that taxi was the only
dataset with rejected rows. All 8,490 rejections were the intentional
\texttt{duplicate\_record} rows. Entries with zero processed rows represent
valid rebuild or product-refresh stages rather than failed ingestion.

\section{Evaluation}

Taxi is the dominant operation because its release contains over half a million records. It processed 517,165 rows in 104.352 seconds, inserted 508,675, and rejected 8,490 duplicates; validation took 20.042 seconds. Weather inserted all 168 rows in 6.581 seconds with 0.751 seconds of validation. Air quality also inserted all 168 rows in 10.701 seconds with 0.881 seconds of validation.


Storage after the final run was 1.63 GB in Bronze, 2.83 GB in Silver, and 2.47 GB in Gold. The update files used 23.73 MB, rejected records used 24.78 MB, and the monitoring table used only 93.91 KB.

The result shows a reasonable operational trade-off. The release files and
monitoring log are small compared with the maintained data layers. Retaining
rejected rows consumes modest space but provides an auditable explanation for
why records were not accepted. The major scalability limitation is taxi
duplicate detection without a source ID. Stable trip identifiers and
partition-aware refreshes would reduce future scans and full-product rebuilds.

\section{Discussion and future direction}

The design keeps the operations layer small by reusing the original loading,
validation, Silver, Gold, and product functions. This means the first release
and later releases use the same definition of a valid record. Insert-only
processing protects Bronze history, while the duplicate check makes a repeated
release safe. Weather and air quality can use their observation keys; taxi has
no stable source identifier, so a SHA-256 hash of shared fields is a practical
way to detect the copied duplicate rows. The final results support this choice:
all 8,490 intentional taxi duplicates were isolated, while all valid weather
and air-quality rows were inserted.

The approach is deliberately simple rather than fully incremental everywhere.
Only products affected by changed sources refresh, but an affected product is
rebuilt in full because these are whole-history aggregates. This is easy to
verify and appropriate for the current dataset. At a larger scale, stable taxi
trip IDs and month-level product refreshes would reduce duplicate-comparison
and recomputation cost. The monitoring table can also support future alerting
when rejection rates or execution times exceed expected levels.
\section{Reproducibility}

From the project root, after the original pipeline has completed:

\begin{verbatim}
python scripts/operations/create_update_release.py
python scripts/operations/run_pipeline.py
python scripts/operations/show_monitoring.py
python scripts/operations/evaluate.py
\end{verbatim}

The release summary is stored at \texttt{data/updates/summary.json}; measured
results are stored at \texttt{data/benchmark/operations\_evaluation.json}.
The accompanying documentation in \texttt{docs/week3/} explains each stage and
its outputs.

\end{document}
